\section{A finite plus construction}\label{sec:plus}

We next recall how a finitely normally generated perfect subgroup
produces a finite $D(2)$ complex. This is the finite form of Quillen's
plus construction~\cite{Quillen71} used by Mannan~\cite[Section~2]{Mannan09}.
We give the construction because its cellular splitting proves both
the required cohomological vanishing and the preservation of the twisted
homology used in Section~\ref{sec:tower}.

Recall that a group $K$ is \emph{perfect} if $K=[K,K]$, or equivalently
if its abelianization $K_{\ab}$ is zero. A normal subgroup of a group
$P$ is \emph{finitely normally generated} if it is the normal closure
in $P$ of finitely many elements. This condition does not require
finite generation as an abstract group.

\begin{proposition}[Finite plus construction]\label{prop:plus}
Let $B$ be a finite connected CW complex of dimension at most two,
with $\pi_1(B)=P$. Suppose that the normal closure $K$ of elements
$t_1,\ldots,t_r\in P$ is perfect, and put $G=P/K$.
There is a finite complex $X\supseteq B$, obtained by attaching
$r$ cells of dimension two and $r$ cells of dimension three, such that:
\begin{enumerate}
\item the inclusion induces the quotient $P\to G$ on fundamental groups;
\item $H_i(\widetilde X;\Z)=0$ for every $i>2$, and
      $H^3(X;M)=0$ for every left $\Z G$-module $M$;
\item for every right $\Z G$-module $N$, the inclusion induces
      $H_*(B;N)\cong H_*(X;N)$, with the coefficients on $B$
      pulled back along $P\to G$.
\end{enumerate}
\end{proposition}

\begin{proof}
Attach two-dimensional cells along loops representing $t_1,\ldots,t_r$,
and call the resulting complex $B'$. The van Kampen theorem gives
$\pi_1(B')=G$. Let $\overline B$ be the full preimage of $B$ in the
universal cover $\widetilde{B'}$. Since $P\to G$ is onto,
$\overline B$ is connected; it is the covering of $B$ corresponding
to $K$. Hence
\[
 H_1(\overline B;\Z)=K_{\ab}=0.
\]

Write $R=\Z G$ and regard cellular chains of these covers as left
$R$-modules. The relative cellular complex
$C_*(\widetilde{B'},\overline B)$ consists of $R^r$ in degree two
and zero in every other degree. The homology exact sequence of the
pair therefore gives a surjection
\begin{equation}\label{eq:relative-surjection}
 H_2(\widetilde{B'};\Z)
 \longrightarrow H_2(\widetilde{B'},\overline B;\Z)=R^r.
\end{equation}
Choose classes $\alpha_1,\ldots,\alpha_r$ mapping to the standard
basis vectors, where the basis records chosen lifts and orientations
of the new cells. The space $\widetilde{B'}$ is simply connected, so
the degree-two Hurewicz theorem represents each $\alpha_j$ by a map
$S^2\to\widetilde{B'}$. Project these maps to $B'$ and attach one
three-dimensional cell along each projected map. This defines $X$.
Only finitely many cells have been added, and attaching the last cells
does not change the fundamental group.

The preimage of $B'$ in $\widetilde X$ is its universal cover, so we
may retain the notation $\widetilde{B'}$ and $\overline B$. In the
chosen cellular bases,
\[
 C_3(\widetilde X)=R^r,
 \qquad C_2(\widetilde X)=C_2(\overline B)\oplus R^r.
\]
Let $q:C_2(\widetilde X)\to R^r$ be projection onto the second
summand. The choice in \eqref{eq:relative-surjection} says exactly that
\begin{equation}\label{eq:split-boundary}
 q\,\partial_3=\id_{R^r}.
\end{equation}
Indeed, the relative cellular boundary of the $j$th three-dimensional
cell is the $j$th basis vector; in the two-dimensional complex
$\widetilde{B'}$, a homology class is represented by its cellular
two-cycle without any ambiguity from three-boundaries.

Equation~\eqref{eq:split-boundary} proves that $\partial_3$ is
injective. There are no cells above dimension three, so
$H_i(\widetilde X;\Z)=0$ for all $i>2$. It also gives the
cohomological assertion directly. For any $R$-linear map
$\varphi:C_3(\widetilde X)\to M$, its extension
$\varphi q:C_2(\widetilde X)\to M$ satisfies
\[
 (\varphi q)\partial_3=\varphi.
\]
Thus the cochain boundary
$\Hom_R(C_2(\widetilde X),M)\to\Hom_R(C_3(\widetilde X),M)$
is onto, and $H^3(X;M)=0$.

Finally, the relative cellular complex of $(\widetilde X,\overline B)$
is
\begin{equation}\label{eq:relative-complex}
 0\longrightarrow R^r\xrightarrow{\,\id\,}R^r
   \longrightarrow0,
\end{equation}
in degrees three and two. It is contractible. The short exact sequence
of absolute and relative cellular complexes is split in each degree,
since the cells of $\overline B$ form subsets of the chosen cellular
bases. Tensoring with any right $R$-module $N$ therefore leaves a short
exact sequence of complexes whose relative term is still contractible.
Its homology exact sequence gives
\[
 H_*(N\otimes_R C_*(\overline B))
   \cong H_*(N\otimes_R C_*(\widetilde X)).
\]
The left side computes $H_*(B;N)$: the coefficient action of $P$
factors through the deck group $G$ of $\overline B\to B$.
The right side computes $H_*(X;N)$, proving the last assertion.
\end{proof}

The splitting in \eqref{eq:split-boundary} is stronger than ordinary
third-homology vanishing: it makes the cohomology calculation work for
every coefficient module. In the application, we will choose $B$ and
$K$ so that a character of $G$ has vanishing second homology on $B$,
yet detects a torsion element of $G$. Proposition~\ref{prop:plus}
preserves that homology vanishing when passing to $X$.
